% from main.tex
%%%%%%%%%%  scan idx 98  |  folio 85  %%%%%%%%%%
Si \(\mathcal C\) et \(\mathcal C'\)
sont des \(S\)-catégories triangulées et si \(F\) est un \(S\)-foncteur exact, alors \(F\) définit un \(S\)-foncteur cartésien \(\Trc(\mathcal C)\longrightarrow \Trc(\mathcal C')\).

\textbf{1.1.4.} Soit \(\mathcal C\) une \(S\)-catégorie triangulée, et soit \(\mathcal C'\) une sous-\(S\)-catégorie fibrée pleine de \(\mathcal C\). On dit que \(\mathcal C'\) est une \emph{sous-\(S\)-catégorie triangulée} de \(\mathcal C\) si, pour tout \(X\in\Ob S\), \(\mathcal C'_X\) est une sous-catégorie triangulée de \(\mathcal C_X\) ([V] chap.\ I \S~1 2.3).

\textbf{1.1.5.} Soit \(\mathcal C\) une \(S\)-catégorie additive. On définit (SGA 4 XVII 2.2, 2.4) une \(S\)-catégorie additive (resp.\ triangulée) \(\Cpx(\mathcal C)\), (resp.\ \(\Kpx(\mathcal C)\)), dont la fibre en \(X\in\Ob S\) est la catégorie \(\Cpx(\mathcal C_X)\) (resp.\ \(\Kpx(\mathcal C_X)\)) des complexes\footnote{à différentielle de degré \(+1\).} (resp.\ complexes «~à homotopie près~») de \(\mathcal C_X\). Si \(\mathcal C\) est une \(S\)-catégorie abélienne plate, on définit (loc.\ cit.) une \(S\)-catégorie triangulée \(\Dpx(\mathcal C)\), dont la fibre en \(X\in\Ob S\) est la catégorie dérivée \(\Dpx(\mathcal C_X)\); si \(f:X\to Y\) est une flèche de \(S\), nous noterons \(f^*:\Dpx(\mathcal C_Y)\longrightarrow \Dpx(\mathcal C_X)\) le foncteur image inverse, qui se déduit de \(f^*:\Kpx(\mathcal C_Y)\longrightarrow \Kpx(\mathcal C_X)\) par simple passage au quotient\footnote{si \(a\in\Ob \Dpx(\mathcal C_Y)\) (resp.\ \(F/\Dpx(\mathcal C_Y)\)), nous écrirons souvent, en abrégé, \(a_X\) au lieu de \(f^*(a)\).}. On définit des sous-\(S\)-catégories triangulées pleines \(\Kpx^*(\mathcal C)\) (resp.\ \(\Dpx^*(\mathcal C)\)) \((*=+,-,b)\) par les conditions habituelles de degré.

\textbf{1.2.} \emph{Terminologie.} Soient \(S\) un site, \(\mathcal C\) une \(S\)-catégorie abélienne plate (1.1), \(\mathcal C_0\) une sous-\(S\)-catégorie fibrée strictement pleine.

Soit \(F\in\Ob\mathcal C_X\), pour \(X\in\Ob S\). L'ensemble des \(f:U\to X\) de \(S_{/X}\) tels qu'il existe un épimorphisme \(E\longrightarrow f^*F\) avec \(E\in\Ob\mathcal C_{0,U}\) est un crible de \(X\). Si celui-ci est couvrant, nous dirons que \(F\) est de \(\mathcal C_0\)-\emph{type fini} (ou simplement de type fini s'il n'y a pas de confusion à redouter).

